\begin{titlepage}
\thispagestyle{empty}
\begin{tikzpicture}[remember picture, overlay]
  \fill[oliveink] (current page.north west) rectangle (current page.south east);
  \fill[papergold!18] ([xshift=-0.7cm,yshift=-2.1cm]current page.north east) circle (6.5cm);
  \fill[white!8] ([xshift=1.3cm,yshift=2.1cm]current page.south west) circle (5.0cm);

  \draw[papergold!80, line width=1.4pt]
    ([xshift=1.1cm,yshift=-1.1cm]current page.north west)
    rectangle
    ([xshift=-1.1cm,yshift=1.1cm]current page.south east);

  \node[
    anchor=north west,
    fill=papergold,
    text=oliveink,
    font=\bfseries\large,
    rounded corners=4pt,
    inner xsep=10pt,
    inner ysep=6pt
  ] at ([xshift=1.6cm,yshift=-1.5cm]current page.north west) {QUYỂN XIV};

  \node[
    anchor=north,
    text=white,
    font=\large\itshape
  ] at ([yshift=-1.7cm]current page.north) {Truyện Tích Cảm Hứng Song Ngữ};

  \node[
    anchor=west,
    text=white,
    font=\fontsize{22}{26}\selectfont\bfseries
  ] at ([xshift=1.8cm,yshift=4.2cm]current page.center) {100 CÂU HỎI};

  \node[
    anchor=west,
    text=papergold,
    font=\fontsize{22}{26}\selectfont\bfseries
  ] at ([xshift=1.8cm,yshift=2.3cm]current page.center) {HỌC TRÒ HAY HỎI};

  \node[
    anchor=west,
    text=white,
    font=\fontsize{22}{26}\selectfont\bfseries
  ] at ([xshift=1.8cm,yshift=0.4cm]current page.center) {VỀ CHUYỆN HỌC};

  \draw[papergold, line width=2pt]
    ([xshift=1.8cm,yshift=-1.0cm]current page.center) --
    ([xshift=10.1cm,yshift=-1.0cm]current page.center);

  \node[
    anchor=west,
    text=white,
    font=\large
  ] at ([xshift=1.8cm,yshift=-2.1cm]current page.center) {Một quyển đi thẳng vào những câu hỏi học sinh thường hỏi về việc học, điểm số và đường đời};

  \node[
    anchor=north west,
    fill=white,
    text=black,
    rounded corners=8pt,
    text width=11.7cm,
    inner sep=14pt,
    align=left,
    font=\normalsize
  ] at ([xshift=1.9cm,yshift=-14.9cm]current page.north west) {%
    \textbf{Tinh thần quyển này:} không trả lời kiểu lên lớp khô cứng. Mỗi câu hỏi được mở ra bằng những băn khoăn rất thật của học sinh về học hành, điểm số, nỗi sợ, động lực và cách sống.\\[6pt]
    \textbf{Cách dùng:} đọc để tự học, học tiếng Anh qua câu ngắn, và dùng như chất liệu giảng dạy kỹ năng học tập cho học sinh.\\[6pt]
    \textbf{Điểm tựa:} ngắn, rõ, gần lớp học, gần tâm lý học sinh, và đi từ chuyện học sang chuyện làm người.
  };

  \node[
    anchor=south,
    text=papergold,
    font=\large\itshape,
    align=center
  ] at ([yshift=2.1cm]current page.south) {%
    ``Có những câu hỏi học sinh hỏi bằng miệng.\\
    Có những câu khác nằm trong cách các em học, sợ, im, và bỏ cuộc.''%
  };

  \node[
    anchor=south,
    text=white!85,
    font=\normalsize,
    align=center
  ] at ([yshift=1.1cm]current page.south) {Dành cho người học, người dạy học, và người muốn nói về việc học theo cách thật hơn};
\end{tikzpicture}
\end{titlepage}
\clearpage
